\documentclass[11pt]{article}
\usepackage[margin=1in]{geometry}
\usepackage{amsmath}
\usepackage{booktabs}
\usepackage{enumitem}
\usepackage{hyperref}

\title{\textbf{ADR-07: Customer Request Topic Profile via LDA Modeling}}
\author{Architecture Decision Record}
\date{\today}

\begin{document}

\maketitle

\section{Status}
\textbf{Accepted}

\section{Context}

Understanding the thematic structure of customer communications requires sophisticated analysis beyond simple categorization. Product teams need insights into underlying customer needs, support teams require understanding of topic complexity, and strategic planners need visibility into emerging themes that may indicate new product opportunities or systemic issues.

\subsection{Business Requirements}
\begin{itemize}
\item Discover latent themes in customer communications for strategic product planning
\item Identify emerging topics that may indicate new customer needs or product gaps
\item Provide product teams with data-driven insights for feature prioritization
\item Enable support teams to understand topic complexity and specialization requirements
\item Support strategic decision-making through thematic analysis of customer communications
\end{itemize}

\subsection{Technical Challenges}
\begin{itemize}
\item Unsupervised learning requires careful parameter tuning for meaningful results
\item Topic interpretability depends on sophisticated text preprocessing and visualization
\item Model validation challenging without labeled ground truth for topic assignments
\item Scalability requirements for larger datasets and real-time topic detection
\item Balance between topic specificity and interpretable generalization
\end{itemize}

\section{Decision}

We implemented \textbf{Latent Dirichlet Allocation (LDA) Topic Modeling} using Gensim for comprehensive thematic analysis:

\subsection{Core Methodology}
\begin{itemize}
\item \textbf{LDA Implementation}: Gensim LdaModel with 5 distinct topics for optimal interpretability
\item \textbf{Advanced Preprocessing}: Multi-stage text cleaning with dictionary filtering for quality
\item \textbf{Bag-of-Words Representation}: Document vectorization optimized for topic modeling
\item \textbf{Parameter Optimization}: 10 passes, chunk size 100, random state 100 for reproducibility
\item \textbf{Dashboard Visualization}: Professional multi-panel display with color-coded topic themes
\end{itemize}

\subsection{Text Processing Pipeline}
\begin{itemize}
\item \textbf{Basic Cleaning}: Lowercase conversion, punctuation removal, stopword filtering
\item \textbf{Length Filtering}: Retain words longer than 2 characters for semantic relevance
\item \textbf{Dictionary Construction}: Gensim Dictionary with extremes filtering for optimal vocabulary
\item \textbf{Corpus Creation}: Bag-of-words transformation for LDA model consumption
\item \textbf{Quality Control}: Remove rare words (<5 documents) and overly common words (>50\% documents)
\end{itemize}

\subsection{Visualization Strategy}
\begin{itemize}
\item \textbf{Dashboard Layout}: Vertical subplot arrangement with 5 topic panels
\item \textbf{Word Importance Display}: Horizontal bar charts showing top 8 words per topic
\item \textbf{Color Coding}: Magma colormap with distinct colors for each topic theme
\item \textbf{Professional Styling}: Clean typography, consistent spacing, and stakeholder-ready presentation
\end{itemize}

\section{Rationale}

\subsection{Why LDA Over Alternative Topic Models}
\begin{itemize}
\item \textbf{Interpretability}: LDA provides probabilistic topic assignments with clear word-topic relationships
\item \textbf{Maturity}: Well-established algorithm with robust Gensim implementation and extensive documentation
\item \textbf{Flexibility}: Configurable topic numbers allowing adaptation to dataset characteristics
\item \textbf{Business Alignment}: Topic-word distributions directly translate to business-understandable themes
\item \textbf{Scalability}: Efficient implementation supporting larger datasets with reasonable computational requirements
\end{itemize}

\subsection{Parameter Selection Rationale}
\begin{itemize}
\item \textbf{5 Topics}: Optimal balance between specificity and interpretability for 50-document dataset
\item \textbf{10 Passes}: Sufficient iterations for model convergence without over-processing
\item \textbf{Dictionary Filtering}: Removes noise while preserving semantic richness
\item \textbf{8 Words per Topic}: Adequate representation for theme understanding without information overload
\end{itemize}

\subsection{Alternative Approaches Rejected}
\begin{itemize}
\item \textbf{K-means Clustering}: Less interpretable results, lacks probabilistic topic assignment
\item \textbf{Manual Theme Identification}: Scale limitations and subjective bias concerns
\item \textbf{Simple Word Frequency}: Missing thematic relationships and contextual understanding
\item \textbf{Neural Topic Models}: Excessive complexity for dataset size and interpretability requirements
\item \textbf{Hierarchical Clustering}: Difficult to determine optimal cluster numbers and less actionable results
\end{itemize}

\section{Consequences}

\subsection{Positive Outcomes}
\begin{itemize}
\item \textbf{Strategic Intelligence}: Clear thematic breakdown enabling product roadmap alignment with customer needs
\item \textbf{Emerging Issue Detection}: Early identification of new customer concerns or product gaps
\item \textbf{Content Strategy}: Data-driven content development aligned with actual customer topic interests
\item \textbf{Team Specialization}: Support team training and specialization based on topic complexity
\item \textbf{Stakeholder Communication}: Executive-ready visualization facilitating strategic discussions
\end{itemize}

\subsection{Technical Benefits}
\begin{itemize}
\item \textbf{Reproducible Analysis}: Fixed random state ensures consistent results for comparative analysis
\item \textbf{Scalable Framework}: Gensim implementation scales to larger datasets efficiently
\item \textbf{Extensible Architecture}: Model easily extended for additional analysis dimensions
\item \textbf{Quality Assurance}: Dictionary filtering ensures high-quality topic generation
\end{itemize}

\subsection{Implementation Trade-offs}
\begin{itemize}
\item \textbf{Parameter Sensitivity}: Topic quality dependent on hyperparameter tuning
\item \textbf{Interpretability Challenges}: Some topics may require subject matter expertise for proper interpretation
\item \textbf{Document Length Dependency}: Performance varies with document length and vocabulary diversity
\item \textbf{Language Specificity}: Current preprocessing optimized for English text only
\end{itemize}

\section{Implementation Details}

\subsection{LDA Model Configuration}
\begin{align}
\text{Documents} &\rightarrow \text{Preprocessing} \rightarrow \text{Clean Tokens} \\
\text{Dictionary} &= \text{corpora.Dictionary}(\text{Clean Tokens}) \\
\text{Filtered Dict} &= \text{filter\_extremes}(\text{no\_below=5, no\_above=0.5}) \\
\text{Corpus} &= [\text{dictionary.doc2bow}(doc) \text{ for } doc \text{ in documents}] \\
\text{LDA Model} &= \text{LdaModel}(\text{corpus, num\_topics=5, passes=10})
\end{align}

\subsection{Visualization Dashboard Parameters}
\begin{itemize}
\item \textbf{Figure Layout}: 5 subplots, 10×15 total size for optimal readability
\item \textbf{Color Scheme}: Magma colormap with 0.0-0.8 range for topic distinction
\item \textbf{Word Display}: Top 8 words per topic with importance weights
\item \textbf{Bar Orientation}: Horizontal layout for word label readability
\item \textbf{Typography}: Consistent font sizing with clear topic labeling
\end{itemize}

\subsection{Performance Characteristics}
\begin{itemize}
\item \textbf{Training Time}: ~30 seconds for 50 documents with 10 passes
\item \textbf{Memory Usage}: Approximately 20MB for complete model and visualization
\item \textbf{Topic Coherence}: Qualitative assessment shows semantically meaningful topic separation
\item \textbf{Scalability}: Linear scaling with document count, quadratic with vocabulary size
\end{itemize}

\section{Success Metrics}

\subsection{Model Quality Indicators}
\begin{itemize}
\item Topics show clear semantic differentiation upon expert review
\item Top words per topic form coherent themes understandable to business stakeholders
\item Model convergence achieved within 10 passes (perplexity stabilization)
\item Topic distribution provides actionable insights for strategic planning
\end{itemize}

\subsection{Business Impact Measures}
\begin{itemize}
\item Product teams incorporate topic insights into feature prioritization discussions
\item Support team specialization decisions informed by topic complexity analysis
\item Strategic planning sessions reference topic trends for market opportunity identification
\item Content development aligned with identified customer topic interests
\end{itemize}

\section{Future Evolution}

\subsection{Model Enhancement Opportunities}
\begin{itemize}
\item \textbf{Dynamic Topic Numbers}: Automatic optimal topic count determination using coherence scoring
\item \textbf{Temporal Topic Modeling}: Track topic evolution over time for trend analysis
\item \textbf{Hierarchical Topics}: Multi-level topic modeling for detailed thematic analysis
\item \textbf{Topic Labeling}: Automatic topic name generation for improved interpretability
\item \textbf{Document Classification}: Assign new documents to existing topic framework
\end{itemize}

\subsection{Integration Roadmap}
\begin{itemize}
\item Real-time topic detection for incoming customer communications
\item Dashboard integration for continuous topic monitoring and trend analysis
\item Machine learning pipeline integration for automated topic-based routing
\item API development for programmatic access to topic analysis capabilities
\end{itemize}

\subsection{Advanced Applications}
\begin{itemize}
\item \textbf{Product Intelligence}: Correlation analysis between topic trends and product usage metrics
\item \textbf{Competitive Analysis}: Topic comparison across different customer segments or time periods
\item \textbf{Content Personalization}: Dynamic content recommendations based on individual topic preferences
\item \textbf{Strategic Forecasting}: Predictive analysis of emerging customer needs through topic trend modeling
\end{itemize}

\end{document}